\documentclass[10pt,a4paper]{amsart}
\usepackage[margin=19mm]{geometry}
\usepackage{amsmath,amssymb,mathtools}
\usepackage[scr=rsfs,cal=euler]{mathalfa}
\usepackage{xfrac}
\usepackage[unicode,hidelinks,bookmarks=false]{hyperref}
\newcommand{\ie}{i.e.\ }
\newcommand{\cf}{cf.\ }
\newcommand{\resp}{respectively\ }
\newcommand{\Subsection}{\subsection}
\providecommand{\nobreakdash}{\nobreak\hbox{-}}
\setlength{\emergencystretch}{2em}
\multlinegap0pt
\usepackage{amssymb}
\newtheorem{theorem}{Theorem}
\title{The construction of $q$-analogues via $_3\phi_2$-series and $q$-difference equations}
\author{John M. Campbell}
\date{Source: arXiv:2602.14314v1 (CC BY 4.0)}
\begin{document}
\maketitle

\noindent\emph{Source and licence.} The construction of $q$-analogues via $_3\phi_2$-series and $q$-difference equations. Version arXiv:2602.14314v1 (CC BY 4.0), distributed by arXiv under the Creative Commons Attribution 4.0 International licence, http://creativecommons.org/licenses/by/4.0/, as recorded in the arXiv licence metadata for this exact version and shown on the abstract page https://arxiv.org/abs/2602.14314v1. Full licence notice: \texttt{licenses/arxiv-2602.14314-CC-BY-4.0.txt}.\par\smallskip

\noindent\textit{Editorial scope.} Counted unit: the article's theorem theoremquarter (source0.tex lines 717--741). The article's complete original proof of it is delivered with it (source0.tex lines 743--843, one \texttt{proof} environment of 101 lines). No mathematics is written, repaired or re-derived: the statement and the proof are the article's own lines, in the article's own notation, and no retained line is substituted or re-flowed.  Retained uncounted as the source-local context the proof needs: lines 207--207; lines 710--714.  Nothing was omitted from any retained span, so the package carries no omission record.  No retained line contains a citation, so the wrapper carries no bibliography and no bibliography entry.  No retained line contains a figure environment, an image-inclusion command or drawing code, so no image is delivered and the package declares no asset dependency.  Editorial formatting only, in addition: an AMS \texttt{amsart} preamble that restates the article's own declarations and macros used by the retained lines, namely the environments \texttt{theorem} on separate counters, together with the standard packages those lines need.  The retained environments print in this document as Theorem 1 in order of appearance, whereas the article prints the same items with the numbering of its own full document.  The complete native source of the article as distributed in the arXiv e-print for this exact version is bundled unchanged as \texttt{sources/194701/source0.tex}.
\subsection{{\tt EKHAD}-normalization}\label{sectionEKHAD}

\begin{equation}\label{prototypeqF}
 F(n, k) := q^{k} \left[ \begin{matrix} 
 q^{a}, q^{b} \vspace{1mm} \\ 
 q^{n+c}, q^{n+d} \end{matrix} \, \Bigg| \, q \right]_{k}, 
\end{equation}

\begin{theorem}\label{theoremquarter}
 Letting

 \ 
 
\noindent $ p(n) = q^{a+b+c+1}-q^{a+b+c+d+n} + q^{a+b+d+1} -  q^{a+c+d+n+1}-q^{b+c+d+n+1}+q^{2 c+2 d+3 n}$, 

 \ 

\noindent the relation 
\begin{multline*}
 (q^{a+b} - q^{c+d}) (q^{a+b+1} - q^{c+d}) 
 {}_{3}\phi_{2}\!\!\left[ \begin{matrix} 
 q^{a}, q^{b}, q \vspace{1mm}\\ 
 q^{c}, q^{d} \end{matrix} \ \Bigg| \ q; q 
 \right] = \\ 
 \sum_{n=0}^{\infty} q^{n} 
 \left[ \begin{matrix} 
 q^{-a+c}, q^{-b+c}, q^{-a+d}, q^{-b+d} \vspace{1mm} \\ 
 q^{c}, q^{d} \end{matrix} \, \Bigg| \, q \right]_{n} 
 \frac{p(n)}{ 
 \big( q^{1-a-b+c+d};q^2 \big)_{n} \big( q^{2-a-b+c+d};q^2 \big)_{n} }
 \end{multline*}
 holds if $\big| \frac{q^{a + b + 1}}{q^{c + d}} \big| < 1 < |q|$. 
\end{theorem} 

\begin{proof}
 For convenience, we write $\alpha = q^a$, $\beta = q^b$, $\gamma = q^c$,  and $\delta = q^d$.  Define  $$ \overline{F}(n, k; q) :=  
 q^{k} \left[ \begin{matrix} 
 \alpha, \beta \vspace{1mm} \\ 
 \gamma q^{n}, \delta q^{n} \end{matrix} \, \Bigg| \, q \right]_{k} 
 \left[ \begin{matrix} 
 \frac{\gamma}{\alpha}, \frac{\gamma}{\beta}, \frac{\delta}{\alpha}, \frac{\delta}{\beta} \vspace{1mm} \\ 
 \gamma, \delta \end{matrix} \, \Bigg| \, q \right]_{n} \frac{1}{ \left( \frac{\gamma \delta}{\alpha \beta q}; q \right)_{2n} },  $$  i.e., as the 
 {\tt EKHAD}-normalized version of \eqref{prototypeqF}. We then let 
\begin{multline*}
 \overline{R}(n, k; q) = \\ q^{n-k} \frac{ -\alpha \beta \gamma \delta q^{k+n} + \alpha \beta \gamma 
 q + \alpha \beta \delta q - \alpha \gamma \delta q^{n+1} - \beta \gamma \delta q^{n+1} 
 + \gamma^{2} \delta^{2} q^{3n+k} }{(\alpha \beta - \gamma \delta q^{2n}) (\alpha \beta q - \gamma \delta q^{2n})}. 
\end{multline*}
 We then define 
 $\overline{G}(n, k; q) := \overline{R}(n, k; q) \overline{F}(n, k; q)$. We may thus verify that 
 the $q$-WZ difference equation 
\begin{equation}\label{firstdiff}
 \overline{F}(n+1, k; q) - \overline{F}(n, k; q) = \overline{G}(n, k+1; q) - \overline{G}(n, k; q) 
\end{equation}
 holds. The telescoping argument from Section \ref{sectionEKHAD}
 applied to \eqref{firstdiff} then 
 gives us that 
\begin{equation}\label{kbarN}
 \sum_{k=0}^{\overline{k} - 1} \left( \overline{F}(N,k;q) - \overline{F}(0,k;q) \right) 
 = \sum_{n=0}^{N-1} \left( \overline{G}(n,\overline{k};q) - \overline{G}(n,0;q) \right). 
\end{equation}
 For $|q| > 1$, we have that 
\begin{multline*}
 \lim_{N \to \infty} \frac{ \overline{F}(N+1,k;q) }{\overline{F}(N,k;q)} 
 = \\
 \lim_{N \to \infty} \frac{ (\gamma q^{N} - \alpha) (\gamma q^{N} - \beta) (\delta q^{N} - \alpha) (\delta q^{N} - 
 \beta) 
 q }{ (1 - \gamma q^{N+k}) (1 - \delta q^{N + k}) (\gamma \delta q^{2N} - \alpha \beta q) (\gamma \delta q^{2N} - 
 \alpha \beta) } = 0, 
\end{multline*} 
 so that $\lim_{N \to \infty} \overline{F}(N,k;q) = 0$. Similarly, and again for $|q| > 1$, we find that 
 $$ \lim_{n \to \infty} \frac{ \overline{G}(n+1,k;q) }{\overline{G}(n,k;q)} 
 = \lim_{n \to \infty} \frac{ \overline{R}(n+1,k;q) }{\overline{R}(n,k;q)} \lim_{n \to \infty} \frac{\overline{F}(n+1,k;q)}{\overline{F}(n,k;q)} = 0, $$
 noting that $ \lim_{n \to \infty} \frac{ \overline{R}(n+1,k;q) }{\overline{R}(n,k;q)} = 1$. 
 Letting $N \to \infty$ in \eqref{kbarN}, we obtain that 
\begin{equation}\label{bothconverge}
 \sum_{k=0}^{\overline{k} - 1} \overline{F}(0, k; q) = \sum_{n=0}^{\infty} \overline{G}(n, 0; q) - \sum_{n = 
  0}^{\infty} \overline{G}(n, \overline{k}; q). 
\end{equation}
 Since $$ \lim_{k \to \infty} \frac{ \overline{F}(0, k+1;q) }{\overline{F}(0, k;q)} = \frac{\alpha \beta q}{\gamma \delta}, $$
 we require the condition $\big| \frac{\alpha \beta q}{\gamma \delta} \big| < 1$
 to ensure the convergence of both sides of \eqref{bothconverge}, so we assume that this condition holds. 
 From the definition of $\overline{R}(n, k;q)$, we find that 
\begin{multline*}
 \left| \overline{R}(n, k; q) \right| = 
 \Bigg| \frac{1}{ \left( \alpha \beta q^{-2n} - \gamma \delta \right) 
 \left( \alpha \beta q^{-2n+1} - \gamma \delta \right)} 
 \big( -\alpha \beta \gamma \delta q^{-2n} + \\ \alpha \beta \gamma q^{-3n - k + 1} 
 + \alpha \beta \delta q^{-3n-k+1} - \alpha \gamma \delta q^{-2n - k + 1} - \beta \gamma \delta q^{-2n-k+1}
 + \gamma^{2} \delta^{2} \big) \Bigg|. 
\end{multline*}
 Consequently, we have that  $$ \left| \overline{R}(n, k; q) \right|   \leq \frac{ |\alpha \beta \gamma \delta| + |\alpha \beta \gamma q|      + 
  |\alpha \beta \delta q| + |\alpha \gamma \delta q| + |\beta \gamma \delta q|     + |\gamma^{2} \delta^{2}| }{ \gamma^{2} \delta^{2}  
  \left| 1 - \frac{\alpha \beta q}{\gamma \delta} q^{-2n-1} \right|    \left| 1 - \frac{\alpha \beta q}{\gamma \delta} q^{-2n} \right| }, $$   and  
    we can conclude that    $ \big| \overline{R}(n, k; q) \big| \leq M$,     for a value $M$ independent of $n$ and $k$,  with 
\begin{equation}\label{boundGwithF}
 \left| \overline{G}(n, k;q) \right| =    \left| \overline{R}(n, k;q) \right| \, \left| \overline{F}(n, k; q) \right|    \leq M \, \left| \overline{F}(n, k;  
  q) \right| 
\end{equation}
 for the same value $M$. Since 
\begin{equation}\label{aktoadd}
 \left| \frac{\overline{F}(n, k+1;q)}{\overline{F}(n,k;q)} \right|    = \left| q \frac{ \left( 1 - \alpha q^{k} \right) \left( 1 - \beta q^{k} 
 \right) }{ \left( 1 - \gamma q^{n+k} \right) \left( 1 - \delta q^{n+k} \right) } \right|, 
\end{equation}
 by writing   $$ \mathcal{A}_{k} = q^{-2n} \frac{\left( 1 - \alpha^{-1} q^{-k} \right) 
 \left( 1 - \beta^{-1} q^{-k} \right) }{ \left( 1 - \gamma^{-1} q^{-n-k} \right) \left( 1 - 
 \delta^{-1} q^{-n-k} \right) }, $$
 this leads us to rewrite \eqref{aktoadd} so that 
\begin{equation*}
 \left| \frac{\overline{F}(n, k+1;q)}{\overline{F}(n,k;q)} \right| 
 = \left| \frac{\alpha \beta q}{\gamma \delta} \mathcal{A}_{k} \right|, 
\end{equation*}
 recalling that $|q| > 1$.  Since $\lim_{k \to \infty} \mathcal{A}_{k} = q^{-2n}$,   we can conclude, again using the condition that $\big|  
 \frac{\alpha \beta q}{\gamma \delta} \big| < 1$,  that 
\begin{equation*}
 \left| \frac{\overline{F}(n, k+1;q)}{\overline{F}(n,k;q)} \right|  < 1. 
\end{equation*}
 Recalling \eqref{boundGwithF}, we have, for $k > \overline{k}_{0}$, that 
\begin{align*}
 \left| \overline{F}(n, k; q) \right| & \leq \rho^{k - \overline{k}_{0}} \left| \overline{F}(n, \overline{k}_{0}; q) \right| \ \text{and} \\ 
 \left| \overline{G}(n, k;q) \right| & \leq \rho^{k - \overline{k}_{0}} M \left| \overline{F}(n, \overline{k}_{0}; q) \right|, 
\end{align*}
 for a positive value $\rho$. 
 So, for $\overline{k} > \overline{k}_{0}$, we find that
\begin{equation}\label{kbarineq}
 \sum_{n=0}^{\infty} \left| \overline{G}(n, \overline{k};q) \right| 
 \leq M \, \rho^{\overline{k} - \overline{k}_{0}} \sum_{n=0}^{\infty} \left| \overline{F}(n, \overline{k}_{0}; q) \right|, 
\end{equation} 
 with the right-hand side of \eqref{kbarineq} approaching $0$
 as $\overline{k} \to \infty$. 
 So, by letting $\overline{k} \to \infty$ in \eqref{bothconverge}, 
 we obtain that 
 $$ \sum_{k=0}^{\infty} \overline{F}(0, k;q) = \sum_{n=0}^{\infty} \overline{G}(n,0;q), $$
 which is equivalent to the desired result. 
\end{proof}

\end{document}
